\documentclass[11pt]{letter}
\usepackage[margin=1in]{geometry}
\usepackage[hidelinks]{hyperref}
\signature{Nuonan Ouyang\\Corresponding Author}
\address{College of Science and Engineering\\James Cook University\\Australia\\nuonan.ouyang@my.jcu.edu.au}
\begin{document}
\begin{letter}{Professor Ahmed Louri\\Editor-in-Chief\\IEEE Transactions on Sustainable Computing}
\opening{Dear Professor Louri,}

Please find enclosed the major revision of manuscript TSUSC-2026-07-0200, ``Shielded Q-learning for Thermally-Safe, Energy-Aware Multi-Model Intrusion Detection Scheduling on Lightweight Edge Devices.''

We have revised the manuscript extensively in response to the Associate Editor and reviewers. The revision retains the submitted FSSQL-R method and operating constraints, but replaces the earlier replay-centered deployment evidence with new hardware experiments. The revised study now includes independent calibration on Raspberry Pi 3B+ and Raspberry Pi 4B 8GB, online closed-loop execution, a conventional threshold controller, matched-seed statistical tests, Q-learning stability analysis, controller and switching overhead, direct POWER-Z KM003C DC input-energy measurements, and an independent one-time evaluation on the official UNSW-NB15 test set using fixed thresholds.

The new evidence also led us to narrow several claims. We no longer present shielding as universally necessary, Q-learning as universally superior, or the guard as an unconditional complete-pipeline safety guarantee. We report adverse results transparently, including the stronger Safe-Greedy result on the Pi 3B+ stationary binding regime and the absence of a consistent physical-energy advantage for FSSQL-R. We believe these revisions materially strengthen the technical validity and deployment relevance of the paper.

A detailed point-by-point response is provided as a separate Summary of Changes document, and supplemental material contains additional calibration, statistical, energy, and official-test details.

Thank you for your consideration of the revised manuscript.

\closing{Sincerely,}
\end{letter}
\end{document}
